% Part: lambda-calculus
% Chapter: introduction
% Section: conversion

\documentclass[../../../include/open-logic-section]{subfiles}

\begin{document}

\olfileid{lam}{int}{con}
\olsection{परिवर्तन आणि न्यूनीकरण}

लॅम्डा कलनात परिवर्तनाचे किंवा न्यूनीकरणाचे तीन प्रकार आहेत.
“रूपांतरण” एकदिश आहे की द्विदिश, यावर परिवर्तन आणि
न्यूनीकरण यांतील फरक अवलंबून असतो; हे आपण पुढे पाहू.

$\alpha$-परिवर्तनात चरांची नावे बदलली जातात.
$\lambda x. x$ आणि $\lambda y.y$ ही दोन्ही पदे एकच फलन
(ओळख फलन) दर्शवतात असे मानणे स्वाभाविक आहे.
ही कल्पना आपण पुढीलप्रमाणे औपचारिक करू:
%READERNOTE{OLINC-332: गोठवलेला स्रोत येथेच थांबतो; आश्वासित औपचारिक व्याख्या किंवा उर्वरित प्रकार या विभागात दिलेले नाहीत.}

\end{document}
